% roadmap.tex — \input dari report.tex (lampiran)
\chapter{Peta Jalan}
\label{ch:roadmap}

\section{Tujuan dan Kontribusi yang Disasar}
Tujuan: laporan yang menegakkan sekaligus menyelesaikan \emph{celah koordinasi} antara
toleransi kesalahan lapis-runtime (OTP/BEAM) dan lapis-orkestrator (Kubernetes).
Kontribusi yang disasar adalah sebuah \emph{protokol koordinasi lintas-lapis} dengan
kontrak kepemilikan/\emph{handoff} antara supervisor OTP dan pengendali Kubernetes.
Cakupan yang disarankan: inti M3 + M7 + M8 ditambah bagian M4 yang khas BEAM; M5
diturunkan ke pembahasan generalisasi; M1/M9 mengutip dan membedakan diri
dari~\cite{roberts2025}.

\section{Progres dan Keputusan}
Pada daftar ini, \cbdone\ berarti selesai dan \cbtodo\ berarti belum.
\begin{itemize}[leftmargin=2.2em,nosep]
  \item[\cbdone] Taksonomi komunikasi A--E (Bab~\ref{ch:background}).
  \item[\cbdone] Lima titik gesekan + skenario + diagram (Bab~\ref{ch:problems}).
  \item[\cbdone] Pertanyaan penelitian RQ1--RQ5 dan mekanisme M1--M9.
  \item[\cbdone] Uji-tekan kebaruan (24 sumber, 15 klaim terkonfirmasi).
  \item[\cbtodo] Keputusan: ambisi cakupan---bundel M3+M7+M8(+M4) vs.\ satu mekanisme tajam.
  \item[\cbtodo] Keputusan: sasaran venue (menentukan format dan ambang rigor).
  \item[\cbtodo] Keputusan: evaluasi \emph{testbed} nyata vs.\ simulasi/model.
\end{itemize}

\section{Fase}

\subsection*{Fase 0 --- \emph{Due diligence} kebaruan (kecil, lebih dulu)}
\begin{itemize}[leftmargin=2.2em,nosep]
  \item[\cbdone] Verifikasi ulang tiga petunjuk yang \emph{abstain} \textbf{(selesai)}:
        \texttt{akkadotnet-healthcheck} \emph{terkonfirmasi} (jembatan keanggotaan$\rightarrow$\emph{probe}
        ada juga di .NET); Beamlens \emph{tanpa} integrasi Kubernetes (alat diagnostik LLM di
        \emph{supervision tree}); Akka Coordination-Lease \emph{terkonfirmasi} memakai CRD
        \texttt{akka.io/v1 Lease} di etcd.
  \item[\cbdone] Telusuri Scholar/DBLP/IEEE~Xplore \textbf{(selesai)}: tidak ditemukan protokol
        koordinasi OTP$\leftrightarrow$K8s yang terbit; karya terbagi dua (FT BEAM internal
        vs.\ FT Kubernetes).
  \item[\cbdone] Pastikan belum ada protokol koordinasi OTP$\leftrightarrow$K8s yang terbit:
        \emph{celah terkonfirmasi}; keyakinan kebaruan naik sedang $\rightarrow$ tinggi untuk
        celah inti (M3/M7/M8).
\end{itemize}

\subsection*{Fase 1 --- Perancangan mekanisme (RQ4)}
\begin{itemize}[leftmargin=2.2em,nosep]
  \item[\cbtodo] Formalkan protokol: mesin keadaan kepemilikan/\emph{handoff}.
  \item[\cbtodo] Definisikan antarmuka OTP$\leftrightarrow$K8s.
  \item[\cbtodo] Tetapkan hukum kendali M3/M7/M8.
  \item[\cbtodo] (opsional) model formal ringan + invarian keselamatan.
\end{itemize}

\subsection*{Fase 2 --- Implementasi acuan + \emph{testbed}}
\begin{itemize}[leftmargin=2.2em,nosep]
  \item[\cbtodo] Aplikasi Phoenix + libcluster + Horde.
  \item[\cbtodo] Komponen koordinasi: operator Bonny + agen sisi BEAM.
  \item[\cbtodo] \emph{Testbed} \texttt{kind}; \emph{build} terkoordinasi dan \emph{baseline}.
\end{itemize}

\subsection*{Fase 3 --- Evaluasi injeksi kesalahan (RQ2/RQ3)}
\begin{itemize}[leftmargin=2.2em,nosep]
  \item[\cbtodo] Matriks kesalahan $\times$ alat $\times$ metrik (Chaos Mesh/LitmusChaos).
  \item[\cbtodo] Metrik: MTTR, kehilangan permintaan, laju kehilangan state, \emph{restart} mubazir.
  \item[\cbtodo] $N$ ulangan + analisis statistik; \emph{baseline} vs.\ terkoordinasi.
\end{itemize}

\subsection*{Fase 4 --- Generalisasi (RQ5)}
\begin{itemize}[leftmargin=2.2em,nosep]
  \item[\cbtodo] Posisikan terhadap Akka~\cite{akkaHealth,akkaLease} dan Orleans~\cite{orleans}.
  \item[\cbtodo] Argumentasikan asimetri lintas-platform.
\end{itemize}

\subsection*{Fase 5 --- Penulisan \& pengiriman}
\begin{itemize}[leftmargin=2.2em,nosep]
  \item[\cbtodo] Kerangka makalah; pustaka; ancaman validitas; paket reproduksibilitas.
\end{itemize}

\section{Sasaran Venue dan Risiko}
Venue Q1 (Scopus) yang sesuai: JSS, EMSE, FGCS, \emph{Software: Practice \& Experience};
atau IEEE TPDS/TCC/TDSC bila mekanismenya cukup teoretis. Risiko utama: keyakinan
kebaruan masih sedang (ditutup Fase~0); beban rekayasa \emph{testbed} (dapat dipersempit
ke satu mekanisme); kerasnya tuntutan rigor evaluasi; serta persepsi ``\emph{niche}'',
yang dijawab oleh bingkai asimetri lintas-platform pada RQ5.
